\section{One diffusion for the whole surface}
\label{sec:lvfield}

The static chapters treated an expiry as a law.  Chapter~2 built the law
first and read the smile off it; Chapter~3 worked in the volatility chart
and policed validity there.  Either way, a fitted surface is at this point a
\emph{family} of laws, one per expiry, related pairwise by the calendar
comparison of \cref{sec:calendar}: at equal log-moneyness, later normalized
calls dominate earlier ones, which is convex order of the normalized
returns.  Convex order is exactly the condition under which the family is
jointly realizable: by Kellerer's theorem \citep{Kellerer1972}, \emph{some}
martingale has these marginals.  But Kellerer's theorem is an existence
statement.  It does not say which martingale, it does not give the process
any structure a computation could use, and it leaves open what additional
information a particular choice smuggles in.

Dupire's observation \citep{Dupire1994} resolves all three questions inside
one class.  Adopt, as a standing hypothesis, the deterministic-carry
setting of \cref{sec:calendar}, under which $Y_\tau=S_\tau/F_{0,\tau}$ is a
single positive martingale whose time-$\tau$ marginal is Chapter~2's
normalized return for the expiry-$\tau$ slice.  Consider diffusions of the
form
\begin{equation}
\dd Y_\tau=\sqrt{v(\tau,Y_\tau)}\;Y_\tau\,\dd W_\tau,\qquad Y_0=1,
\label{eq:lvsde}
\end{equation}
where $v(\tau,y)\ge0$ is a deterministic field, the \emph{local
variance}; its square root $\sigma_{\rm loc}=\sqrt v$ is the local
volatility.  Within
this class the correspondence between coefficient and marginals is exact in
both directions: the field determines every marginal law by solving a
forward equation, and the marginals determine the field by a pointwise
formula, wherever the density is positive.  A valid surface is therefore
not merely \emph{consistent with} some model --- it \emph{is} the
coefficient field of one canonical model, written in different coordinates.
This chapter is about that change of coordinates: how to compute it, why
the obvious way of computing it fails on real data, and what a market's
finitely many quotes do and do not say about the field.

Two remarks bound the object's ambitions before we start.  First, the
diffusion \eqref{eq:lvsde} is a \emph{representative}, not a discovery.  By
Gy\"ongy's theorem \citep{Gyongy1986}, any It\^o martingale (under mild
regularity) with the same marginals is mimicked by \eqref{eq:lvsde} with
$v(\tau,y)$ equal to the conditional expectation of its instantaneous
variance given $Y_\tau=y$: the local-volatility model is the Markovian
projection common to every dynamics the statics allow.  Statements that
depend only on marginal laws --- every European price, every quantity of
Chapters~2--3 --- are safe; statements about joint laws, such as how the
smile moves when the spot moves, are not determined by the surface at all
and are postponed to Chapter~9.  Second, the class \eqref{eq:lvsde}
contains no jumps.  A genuine discontinuity --- an earnings gap, a
scheduled announcement --- can only be imitated by a burst of diffusive
variance; scheduled events belong to a change of clock, which is
Chapter~8's subject.  Throughout this chapter $\tau$ is the
variance time used since Chapter~2, $y=e^{k}$ the normalized strike, and
$c(\tau,y)$ the normalized call: rates, borrow and dividends live inside
the forward (Chapter~6), and the field $v$ is variance per unit variance
time.

The plan follows the two directions of the correspondence.
\Cref{sec:lveq} derives the forward equation, shows that positivity of the
field propagates both static no-arbitrage families at once, and identifies
the inverse formula as a ratio of two objects the reader already owns.
\Cref{sec:lvillposed} demonstrates why that formula must not be applied to
data.  \Cref{sec:lvsheet,sec:lvlattice} build the objects that replace it:
a finite-dimensional field whose validity is a box constraint, and a
discretization that inherits the continuous theorem instead of merely
approximating it.  \Cref{sec:lvinfo} asks what the quotes actually
determine, and \cref{sec:lvcalib,sec:lvexamples} assemble the calibration
and run it on the book's frozen data.
